\section{Supplementary methods}

\subsection{Training data}
The training data come from four sources, all mapped to GTDB R\GTDBRelease \citep{parks2022}.
\begin{itemize}
\item \textbf{Species with a processed reference genome} (\NOriginal rows). Species from the growth-rate table distributed
  with gRodon \citep{weissman2021}, which derives from the compilation of \citet{madin2020}, for which we downloaded a
  genome assembly from NCBI and computed every feature ourselves.
\item \textbf{Madin extension} (\NMadinExtension rows). Species of the same table that had no exact match on the reference
  tree in an earlier version of the pipeline. Genome size, GC content and tRNA count are taken from Madin's species
  table; the gene features come from the genome of the accession Madin's accession table gives for the species.
\item \textbf{Madin genome extension} (\NGenomeExtension rows). Species of the same table with a genome accession but no
  traits in Madin's table. Every feature is computed from that genome: size and GC exactly, tRNA genes with tRNAscan-SE
  \citep{chan2021} using the bacterial or archaeal model, and the gene features as for all genomes. Names such as
  ``\textit{Dictyoglomus} sp.'' are skipped, because the value belongs to an unnamed strain. When the RefSeq record of an
  assembly is suppressed, the GenBank copy is used.
\item \textbf{Literature strains} (\NLiterature rows). Strains from lineages the compilation lacks, each with a doubling
  time read from the primary paper and paired with the genome of the measured strain
  \citep{konneke2005,stieglmeier2014,palatinszky2015,nowka2015,mohammadi2017,lehtovirtamorley2014}. Values that could not
  be confirmed in the primary source, came from enrichment cultures, or described something other than population growth
  (for example a budding cycle) were not used.
\end{itemize}
Each row is placed in GTDB by its genome accession where possible, else by species name, else by genus name. Rows are
grouped into species groups: two rows are in the same group when they share a species name or a GTDB species, followed
transitively, so that a renamed species and the same name placed in two GTDB species stay together in every split.
Table~\ref{tab:literature} lists the literature strains.

\begin{table}[ht]\centering\small
\caption{Doubling times added from the primary literature for nitrifying and methanotrophic strains with a sequenced genome, with the temperature and conditions of the measurement.}\label{tab:literature}
\begin{tabular}{p{3.6cm}rrp{4.8cm}p{4.2cm}}
\toprule
Strain & Doubling (h) & \textdegree C & How measured & Source \\
\midrule
\textit{Nitrosopumilus maritimus} SCM1 & 21 & 28 & minimum generation time, 500 uM ammonium, defined medium & Konneke et al. 2005 Nature 437:543 \\
\textit{Nitrososphaera viennensis} EN76 & 27.5 & 42 & generation time from cell counts, 1 mM pyruvate, 2 mM NH4Cl & Stieglmeier et al. 2014 IJSEM 64:2738 \\
\textit{Nitrososphaera gargensis} Ga9.2 & 113.4 & 46 & mean generation time from total protein, 0.5 mM ammonium & Palatinszky et al. 2015 Nature 524:105 \\
\textit{Nitrospira defluvii} A17 & 37 & 28 & generation time, pure culture (Table 1) & Nowka et al. 2015 Appl Environ Microbiol 81:745 \\
\textit{Nitrospira moscoviensis} NSP M-1 & 32 & 37 & generation time, pure culture (Table 1) & Nowka et al. 2015 Appl Environ Microbiol 81:745 \\
\textit{Nitrobacter hamburgensis} X14 & 43 & 28 & generation time, pure culture (Table 1) & Nowka et al. 2015 Appl Environ Microbiol 81:745 \\
\textit{Methylacidiphilum fumariolicum} SolV & 9.9 & 55 & batch on methane with ammonium, mu 0.07 per h, pH 2.7 & Mohammadi et al. 2017 ISME J 11:945 \\
\textit{Nitrosotalea devanaterra} Nd1 & 57.4 & 25 & mu 0.29 per day at the optimum temperature, pH 5 & Lehtovirta-Morley et al. 2014 FEMS Microbiol Ecol 89:542 \\
\bottomrule
\end{tabular}
\end{table}


\subsection{Label audit}
\paragraph{Species-level screen.} A label is flagged when it is at least \ScreenFoldWord-fold away from both its prediction with the
species held out and the median label of at least two other labeled species in its genus. The source records of every
flagged label are then read in the per-source files of the compilation, which give each record with its temperature and
reference.

\paragraph{Study-level check.} Each training label is traced to the record it came from by matching the value (within
\TraceTolPct\%). Each traced label is compared with evidence that does not come from the same study: the median of other studies'
values for the same species, or, failing that, the median label of at least two relatives in the same genus (or family)
whose labels come from other studies. A study is flagged when at least two of its labels can be compared, their median
ratio to the evidence is at least \StudyFoldWord-fold in either direction, and at least 80\% point the same way. Labels that match no
record but are at least \StudyFoldWord-fold from the closest record are listed separately, since a unit or conversion step leaves
that signature.

\paragraph{Decisions.} A flag is not a removal. A label was removed only when a second source for the same species, or the
study itself, showed that the value is not a maximal growth rate at the stated temperature. Flags with a biological
explanation were kept, for example rich-medium rates of haloarchaea at their optimum, real differences between species of
one genus, and slow species such as \textit{Mycobacterium leprae}. Table~\ref{tab:exclusions} lists the \NExcludedWord
removed values from \NExcludedSourcesWord sources.

\begin{longtable}{p{3cm}p{1.2cm}p{6cm}p{4.4cm}}
\caption{Doubling times removed from the training data because they are not maximal growth rates, with the evidence.}\label{tab:exclusions}\\
\toprule
Species & Label (h) & Evidence & Source \\
\midrule
\endfirsthead
\caption[]{(continued)}\\
\toprule
Species & Label (h) & Evidence & Source \\
\midrule
\endhead
\textit{Methylocystis parvus} & 152 & source gives about 3700 h at 30 C for a methanotroph; same study gives Methylococcus capsulatus about 4600 h against 1.87 h elsewhere & Kevbrina et al. 2001 Microbiology 70:384 (growth of methanotrophs at low temperatures) via Corkrey et al. \\
\textit{Methylomonas methanica} & 112 & source gives about 2700 h; same study as above & Kevbrina et al. 2001 Microbiology 70:384 via Corkrey et al. \\
\textit{Methylosinus sporium} & 187 & source gives about 4600 h; same study as above & Kevbrina et al. 2001 Microbiology 70:384 via Corkrey et al. \\
\textit{Methylosinus trichosporium} & 131 & source gives about 3200 h; same study as above & Kevbrina et al. 2001 Microbiology 70:384 via Corkrey et al. \\
\textit{Flavobacterium segetis} & 433 & recorded at 12.9 C; Bergey's Manual gives 6.8 h; genus relatives median 4.9 h & Yi and Chun 2006 (species description) via Corkrey et al. \\
\textit{Flavobacterium weaverense} & 178 & recorded at 15.6 C; Bergey's Manual gives 2.9 h & Yi and Chun 2006 (species description) via Corkrey et al. \\
\textit{Acidithiobacillus ferrooxidans} & 386 & about 70 times its genus relatives (median 5.5 h); other organisms from the same study look plausible, cause unknown & Franzmann et al. (temperature and iron/sulfur oxidation rates) via Corkrey et al. \\
\textit{Acidithiobacillus thiooxidans} & 74 & about 13 times its genus relatives; same study as above & Franzmann et al. via Corkrey et al. \\
\textit{Klebsiella oxytoca} & 11.7 & from a lag-phase study in meat; genus relatives median 0.3 h & Mellefont (meat industry predictive modelling) via Corkrey et al. \\
\textit{Streptomyces coelicolor} & 0.5 & another source gives 2.04 h; MediaDB values are usually about 1.6 times slower than other sources (defined media), this one is 4 times faster, and 0.5 h is implausible for a filamentous Streptomyces & MediaDB (Richards et al. 2014) via Madin et al.; found by screen\_label\_sources.py \\
\bottomrule
\end{longtable}


\subsection{Features}
\paragraph{Codon usage.} Codon counts of each gene are compared with the codon usage of all genes of the genome by the MILC
distance \citep{supek2005}; the codon usage bias of a genome is the median distance over its ribosomal protein genes. The
released gRodon statistics are computed by \texttt{predictGrowth} in full mode without a temperature, on the same genes,
with the ribosomal protein genes marked as highly expressed; we keep \texttt{CUBHE}, \texttt{ConsistencyHE}, \texttt{CPB}
and the logarithm of gRodon's doubling time.

\paragraph{Lifestyle inputs.} Presence inputs are 1 when at least one protein hits the family (or families) above the
cutoff: aerobic respiration (a heme-copper or cytochrome \textit{bd} oxidase), complex I, methanogenesis (methyl-coenzyme
M reductase), sulfate reduction (DsrD), nitrogen fixation (NifH and at least two NifD/K members beyond those of chlorophyll
synthesis: two in cyanobacteria, which have both reaction centre types, four in other phototrophs), RuBisCO, ammonia or
methane monooxygenase (subunit A, or subunit C by E-value), phototrophy (a reaction centre or photosystem I), flagellar
motility and catalase. Count inputs are the number of terminal oxidases and of hydrogenases. Density inputs are the
numbers of transporter and of regulator genes per thousand genes (Table~\ref{tab:pfam}).

\paragraph{Energy inputs.} For the low-energy test only, \NEnergyFeatures further inputs describe how a genome makes
energy: nitrate respiration (NarG, by its N-terminal domain or domain V, or NapB), denitrification (NosZ), DMSO or TMAO
respiration, fumarate respiration (the \textit{E.~coli} type only), metal reduction outside the cell, the Wood--Ljungdahl
pathway (acetyl-CoA synthase), the number of electron acceptor classes, F-type and V/A-type ATP synthase, Na$^+$-NQR
(the NqrA barrel together with an NqrA-specific domain or NqrM), the Rnf complex, sodium-pumping decarboxylases, the number
of sodium pumps, and fermentation only (fermentation genes and none of the acceptors). Four families were tightened or
dropped after hitting organisms known to lack the function. Nitrite oxidizers are counted as nitrate respirers, because their
nitrite oxidoreductase is a NarG relative. Syntrophy was not used: a rule fixed from the literature before evaluation found
\SynFoundWord of \SynKnownWord known syntrophs and called \SynFalse other genomes, mostly fermenters that carry the same genes.

\begin{longtable}{p{1.6cm}p{8.2cm}p{1.8cm}p{3cm}}
\caption{Pfam families behind the lifestyle inputs of the model and the energy inputs of the low-energy test, with the search cutoff. Accessions are written without version.}\label{tab:pfam}\\
\toprule
Pfam & What it marks & Cutoff & Used in \\
\midrule
\endfirsthead
\caption[]{(continued)}\\
\toprule
Pfam & What it marks & Cutoff & Used in \\
\midrule
\endhead
PF00005 & ABC transporter ATPase & gathering & model input \\
PF00016 & RuBisCO large subunit & gathering & model input \\
PF00072 & response regulator receiver & gathering & model input \\
PF00083 & sugar transporter & gathering & model input \\
PF00115 & heme-copper oxidase subunit I & gathering & model input, energy test \\
PF00119 & F-type ATP synthase subunit a & gathering & energy test \\
PF00124 & photosynthetic reaction centre & gathering & model input, energy test \\
PF00142 & nitrogenase iron protein NifH & gathering & model input \\
PF00148 & nitrogenase MoFe protein NifD/K & gathering & model input \\
PF00196 & LuxR family regulator & gathering & model input \\
PF00199 & catalase & gathering & model input \\
PF00213 & F-type ATP synthase delta (OSCP) & gathering & energy test \\
PF00223 & photosystem I PsaA/PsaB & gathering & model input, energy test \\
PF00361 & complex I membrane subunit & gathering & model input \\
PF00374 & NiFe hydrogenase large subunit & gathering & model input \\
PF00392 & GntR family regulator & gathering & model input \\
PF00440 & TetR family regulator & gathering & model input \\
PF00512 & histidine kinase dimerisation & gathering & model input \\
PF00528 & ABC transporter permease & gathering & model input \\
PF00593 & TonB-dependent receptor & gathering & model input \\
PF00669 & flagellin N-terminal & gathering & model input \\
PF00871 & acetate kinase & gathering & energy test \\
PF01047 & MarR family regulator & gathering & model input \\
PF01496 & V/A-type ATP synthase subunit I & gathering & energy test \\
PF01515 & phosphotransacetylase & gathering & energy test \\
PF01654 & cytochrome bd oxidase subunit I & gathering & model input, energy test \\
PF01991 & V/A-type ATP synthase subunit E & gathering & energy test \\
PF02239 & cd1 nitrite reductase NirS & gathering & energy test \\
PF02241 & methyl-coenzyme M reductase beta & gathering & model input, energy test \\
PF02249 & methyl-coenzyme M reductase alpha & gathering & model input, energy test \\
PF02300 & fumarate reductase subunit C & gathering & energy test \\
PF02313 & fumarate reductase subunit D & gathering & energy test \\
PF02378 & PTS sugar transporter EIIC & gathering & model input \\
PF02461 & ammonia/methane monooxygenase subunit A & gathering & model input \\
PF02665 & respiratory nitrate reductase gamma subunit NarI & gathering & energy test \\
PF02745 & methyl-coenzyme M reductase alpha, N-terminal & gathering & model input, energy test \\
PF02901 & pyruvate formate-lyase & gathering & energy test \\
PF02906 & FeFe hydrogenase & gathering & model input \\
PF03466 & LysR substrate binding & gathering & model input \\
PF03598 & CO dehydrogenase/acetyl-CoA synthase beta & gathering & energy test \\
PF03599 & CO dehydrogenase/acetyl-CoA synthase delta & gathering & energy test \\
PF03892 & periplasmic nitrate reductase NapB & gathering & energy test \\
PF03977 & Na+-pumping decarboxylase beta subunit & gathering & energy test \\
PF04208 & Na+-pumping methyltransferase MtrA & gathering & energy test \\
PF04400 & Na+-NQR maturation protein NqrM & gathering & energy test \\
PF04542 & sigma-70 region 2 & gathering & model input \\
PF04896 & ammonia/methane monooxygenase subunit C & E $<$ 1e-05 & model input \\
PF04976 & DMSO/TMAO reductase anchor DmsC & gathering & energy test \\
PF05896 & Na+-translocating NADH:quinone oxidoreductase NqrA & gathering & energy test \\
PF06808 & TRAP transporter large permease & gathering & model input \\
PF07690 & major facilitator superfamily & gathering & model input \\
PF08679 & dissimilatory sulfite reductase DsrD & gathering & model input, energy test \\
PF11854 & outer membrane porin MtrB/PioB & gathering & energy test \\
PF12833 & AraC family regulator & gathering & model input \\
PF13375 & Rnf complex subunit RnfC & gathering & energy test \\
PF14710 & respiratory nitrate reductase NarG, N-terminal & gathering & energy test \\
PF18764 & nitrous oxide reductase NosZ & gathering & energy test \\
PF22111 & decaheme cytochrome MtrC/MtrF & gathering & energy test \\
PF22112 & decaheme cytochrome OmcA & gathering & energy test \\
PF24836 & NqrA second domain (not in RnfC) & gathering & energy test \\
PF28599 & respiratory nitrate reductase NarG, domain V & gathering & energy test \\
\bottomrule
\end{longtable}


\paragraph{Phylogeny and 16S.} Landmark tips are chosen by farthest-point sampling on each GTDB tree, starting from a tip
drawn with a fixed seed; the vector of a taxon is the mean distance vector of its tips. The 16S reference sequences are
chosen in the same way among the labeled species, and a 16S input is the cosine distance between \KmerK-mer count
profiles. Missing inputs of a labeled row are filled with the median of the bacterial rows; models for genomes that lack a
whole block (the tree position or the 16S gene) are trained without that block.

\paragraph{Input drift check.} Before predicting new genomes, the mean of each input over the new genomes is compared with
the training mean in units of the training standard deviation, separately for bacteria and archaea, and the run stops if
any phylogeny or 16S input has moved by more than \DriftSD standard deviations. This check caught an early version of the
pipeline in which the tree and 16S inputs of the GEM genomes had been computed in a different coordinate system.

\subsection{Model settings}
\begin{table}[ht]\centering\small
\caption{Settings of the two members of the deployed model for each target, chosen by \texttt{tune\_deployed\_model.py}. Settings not listed are scikit-learn defaults; the random forest has \NForestTrees trees.}\label{tab:settings}
\begin{tabular}{llp{9.5cm}}
\toprule
Target & Member & Settings \\
\midrule
own temperature & histogram boosting & min\_samples\_leaf = 20, max\_leaf\_nodes = 7, max\_iter = 200, max\_features = 1.0, learning\_rate = 0.05, l2\_regularization = 5.0 \\
own temperature & random forest & min\_samples\_leaf = 5, max\_samples = 0.7, max\_features = 0.5 \\
20 C & histogram boosting & min\_samples\_leaf = 5, max\_leaf\_nodes = 7, max\_iter = 400, max\_features = 1.0, learning\_rate = 0.05, l2\_regularization = 0.0 \\
20 C & random forest & min\_samples\_leaf = 5, max\_samples = 0.7, max\_features = 0.2 \\
growth temperature & histogram boosting & min\_samples\_leaf = 20, max\_leaf\_nodes = 7, max\_iter = 200, max\_features = 1.0, learning\_rate = 0.05, l2\_regularization = 5.0 \\
growth temperature & random forest & min\_samples\_leaf = 10, max\_samples = None, max\_features = 0.5 \\
\bottomrule
\end{tabular}
\end{table}


\subsection{Comparison with released gRodon and Phydon}
Gene files were written in the layout Phydon expects (nucleotide genes, the ribosomal protein genes named ``ribosomal
protein'', and a list of gene identifiers), so gRodon, Phydon and our model see the same genes. Genome names were given as
GTDB accessions so that Phydon finds each species on its own tree (GTDB R220). For each fold, the species representatives
of the fold's genomes were removed from Phydon's built-in training table before the call, which makes both its
phylogenetic estimate and its distance to the nearest training relative leave the test species out. Phydon was run in its
default geometric-mean mode with gRodon in full mode. Two advantages of the released tools remain: gRodon's regression and
Phydon's weighting model were fitted on data that include our species, and Phydon's training table contains species that
ours does not. The results are in Table~\ref{tab:released}.

\subsection{Phylogenetic tests}
Each species group is placed at the tip of its GTDB species representative, or, when the group is placed only to genus, at
an unused tip of the same genus. The covariance matrix has one block per domain, each the shared root-to-ancestor path
length of the tips, scaled to a mean diagonal of one. Pagel's $\lambda$ multiplies the off-diagonal elements and is
estimated by maximum likelihood in $[0,1]$. The model for slow growth is a logistic regression on the input, growth
temperature and domain, with the leading eigenvectors of the covariance matrix as covariates; we use \KEigen and report fewer and more
as a check, because with half as many the single slow clade of Anaerolineales was not absorbed and an input it happened to carry looked
strongly linked to slow growth.

\subsection{Software}
\begin{table}[ht]\centering\small
\caption{Software versions used for the results in this paper.}\label{tab:software}
\begin{tabular}{ll}
\toprule
Software & Version \\
\midrule
Python & 3.11.14 \\
scikit-learn & 1.9.0 \\
numpy & 2.4.6 \\
pandas & 3.0.5 \\
scipy & 1.17.1 \\
statsmodels & 0.14.6 \\
pyrodigal & 3.7.1 \\
dendropy & 5.0.13 \\
biopython & 1.88 \\
HMMER (hmmsearch) & 3.4 \\
tRNAscan-SE & 2.0.13 \\
barrnap & 0.9 \\
gRodon & 2.7.3 \\
Phydon & 1.1.0 \\
picante & 1.8.4 \\
castor & 1.8.7 \\
ape & 5.8.1 \\
Biostrings & 2.80.2 \\
coRdon & 1.30.0 \\
R & 4.6.1 \\
\bottomrule
\end{tabular}
\end{table}

